\documentclass[landscape,twocolumn,leqno,a4paper,10pt]{report} %leqno affiche les numéros d'équations à gauche

\usepackage[a4paper,top=0.4cm,bottom=1cm,left=0.4cm,right=0.4cm,marginparwidth=0.5cm,footskip=0.6cm]{geometry}



%----------------- Codage de la langue et de la police--------------------------------------------------------------------------
\usepackage[french]{babel} 		% Permet de spécifier au compilateur que l'on désire taper son document en français.
\usepackage[utf8x]{inputenc} 	% Permet de taper directement à l'écran, dans le fichier .tex, les caractères accentués.
\usepackage[T1]{fontenc} 		% Permet juste d'utiliser une nouvelle norme LATEX concernant le codage des caractères.
%\usepackage{lmodern}		% Pack de police --> Pas besoin puisque j'utilise Fourier

%------------------------ Packages utiles-------------------------------------------------------------------------------------------

\usepackage{graphicx}	% Pour insérer des images
\usepackage{caption}		% Légende

\newcommand\hmmax{0}		% Alors c'est une astuce assez mystérieuse pour dépasser la limite de polices latex.
							% En effet, le package bm est lourd et avec tout ce que l'on met, on obtient un message d'erreur : Too many math alphabets used in version normal.			
\newcommand\bmmax{0}		% Cette astuce règle le problème, il faut tout de même mettre ces compteurs avant le package bm

\usepackage{amsmath,amssymb}	% Permettent d'obtenir des caractères mathématiques (de toute beauté !).
\usepackage{amsthm}		% Pour pouvoir faire les newthm
\usepackage{mathtools}	% Corrige des bugs dans amsmath et le complète
\usepackage{numprint}	% Pour les espacements dans les nombres : \numprint{10000} = 10 000
\usepackage{fourier}		% REMARQUE : [upright] impose aux majuscules d'être droite dans le mode mathématique 
\usepackage{pifont} 		% Importe une large gamme de symboles, on utilise \ding{numéro}
\usepackage{xargs} 		% Permet de définir des macros à plusieurs paramètres optionnels.
\usepackage{fancyhdr} 	% Package pour les entêtes et bas de page
\pagestyle{fancy} 			% Accompagne \usepackage{fancyhdr}
\usepackage{bm} 		% Permet de mettre en gras du contenu mathématiques via \bm{..}
\usepackage{mathrsfs} 	% Pour les lettres cursives
\usepackage{enumitem} 	% Pour pouvoir personnaliser les listes

\frenchbsetup{StandardLists=true} % Afin d'éviter tout conflit entre les extensions enumitem et french babel. 

\usepackage{sectsty}	% Pour colorer les titres
\usepackage{titlesec} 	% Pour personnaliser les titres
\usepackage{setspace} 	% Pour changer la taille des interlignes
       	               
\usepackage{array,multirow,makecell}
\usepackage{hlist}		% Package Christian Tellechea sur les listes numérotées gauche vers droite
\setdefaulthlist{left margin = 2mm, label sep= 0.2mm}	% réglages package hlist
				

%------------ Les packages de dessin sous pstricks : il faut changer la compilation de latex pour compiler ces dessins------------------
% Options\Configurer Texmaker\Compil rapide\ Latex+dvips+ps2pdf+Voir PDF

\usepackage[pdf]{pstricks} 	% Pour faire des dessins
\usepackage{pstricks-add,pst-plot,pst-tree,pst-node,pst-math,pst-eucl,pst-func,pst-3dplot}% ,pst-text,pst-eps,pst-fill,pst-grad
\usepackage{multido}		% Pour faire des multido
\usepackage{lastpage} 	% Pour le compteur de la dernière page, on appelle alors celle-ci avec lastage



\usepackage[colorlinks=true, 
			%allcolors=blue
			]{hyperref}		% ATTENTION : doit être importé après la majorité des packages sauf tabularx et breakrules (par exemple)
\hypersetup{			% Paramétrage des hyperliens
	colorlinks = true,   		% Colorise les liens
	breaklinks = true,		% Permet les retours à la ligne pour les liens trop longs
	urlcolor = blue,		% Couleur des hyperliens
	linkcolor = black,		% Couleur des liens internes aux documents (index, figures, tableaux, équations, ...)
	citecolor = blue		% Couleur des liens vers les références bibliographiques
	}




%%%%%%%%%%%%%%%%%%%%%%%%%%%%%%%%%%%%%%%%%%%%%%%%%%
%------------------------- Définition des commandes -----------------------------------------------------------------------------
%%%%%%%%%%%%%%%%%%%%%%%%%%%%%%%%%%%%%%%%%%%%%%%%%%


%--------------------- Raccourcis fontes--------------------------------------------------------------
% REMARQUE : le package fourier met la même police pour mathcal et mathscr, on remet donc mathcal normalement et on utilise la définition de fourier (un peu plus grasse) pour mathscr
\DeclareMathAlphabet{\mathcal}{OMS}{cmsy}{m}{n}	
\DeclareMathAlphabet{\mathscr}{FMS}{futm}{m}{n}

%---------------------------- Raccourcis des espaces--------------------------------------------------------------------------------------------
\newcommand\n{\noindent}
\newcommand\q{\quad}
\newcommand\qq{\quad\quad}

%--------------------- Raccourcis en mode mathématiques-------------------------------------------------------------------------------------
\newcommand{\et}{\text{ et }}
\newcommand{\ou}{\text{ ou }}
\newcommand{\SSI}{\Longleftrightarrow}
\let\si\relax
\newcommand{\si}{\text{ si }}
\newcommand{\sinon}{\text{ sinon }}
\renewcommand{\epsilon}{\varepsilon}		% epsilon beaucoup plus joli
\renewcommand{\geq}{\geqslant} 
\renewcommand{\leq}{\leqslant}
\newcommand{\e}{\mathrm{e}}
\renewcommand{\textbf}[1]{\begingroup\bfseries\mathversion{bold}#1\endgroup}		% Très important, permet de passer en gras même avec le mode mathématiques

\newcommand{\Mod}[1]{\allowbreak\if@display\mkern5mu\else\mkern5mu\fi\left[#1\right]}	% Pour modulo entre crochets [.] avec un espacement convenable 


%---------------------------------- Raccourcis calcul diff----------------------------------------------------------------------------------------
\renewcommand{\d}{\;\mathrm{d}}	% Pour le d droit de l'intégrale car ce n'est pas une variable...


%-------------------------------- Raccourcis des ensembles----------------------------------------------------------------------------
\newcommand\R{\mathbb{R}}
\newcommand\C{\mathbb{C}}
\newcommand\K{\mathbb{K}}
\newcommand\Q{\mathbb{Q}}
\newcommand\D{\mathbb{D}}
\newcommand\Z{\mathbb{Z}}
\newcommand\N{\mathbb{N}}
\newcommand\U{\mathbb{U}}


%------------------------------- Dimensions intelignes, items...-----------------------------------------------------------------------------------------------
\onehalfspacing	% Pour que l'interligne soit de 1.5. Pour que l'interligne soit de la valeur que l'on souhaite, de base : \setstretch{1}

\setlist[1]{font=\bfseries,itemsep=0.2mm,leftmargin=6mm}
\setlist[2]{font=\bfseries,itemsep=0.2mm,leftmargin=4mm}
\setenumerate[1]{font=\bfseries,itemsep=0.2mm,leftmargin=6mm}
\setenumerate[2]{font=\bfseries, label=\alph*., leftmargin=4mm}
%%\setenumerate[profondeur]{font=\bfseries, label=\ding{221} \color{blue}} permet de changer en a) b) les sous-items d'une liste enumerate
%% itemsep = correspond à l'interligne entre chaque item
% de base : labelsep=2mm et permet de changer la distance entre l’étiquette et le texte des éléments
% leftmargin =4mm spécifie la marge à gauche de la liste
\AtBeginDocument{\renewcommand{\labelitemi}{\textbullet}}	% Pour que les puces par défaut soit des points.
%\AtBeginDocument{\renewcommand{\labelitemi}{\ding{86}}}	% Pour que les puces par défaut soit des étoiles.
\AtBeginDocument{\renewcommand{\labelitemii}{$\ast$}}		% Pour que les puces par défaut pour une sous-liste soit des astérisques.
\newcommand{\itemmm}{\item[$\ast$]}				% \itemmm devient l'item \ast
\newcommand{\itemmmm}{\item[\footnotesize$\circledcirc$\normalsize]}
\newcommand{\itemmmmm}{\item[\footnotesize$\astrosun$\normalsize]}
\newcommand{\itemmmmmm}{\item[\footnotesize$\multimap$\normalsize]}
\newcommand{\itemmmmmmmmm}{\item[$\divideontimes$]}
\newcommand{\etoile}{\item[$\star$]}
\newcommand{\etoilee}{\item[$\bigstar$]}
\newcommand{\losange}{\item[$\diamondsuit$]}
\newcommand{\carre}{\item[\LARGE$\centerdot$]}
\newcommand{\carreombre}{\item[\ding{112}]}

\setlength{\columnsep}{30pt} 		% Permet d'écarter les deux colonnes du twocolumn de 30pt, pratique pour écarter le texte de la ligne centrale
\setlength{\columnseprule}{0.5pt}	% Affiche une ligne verticale entre chaque item et entre chaque page du twocolumn
						% \columnseprule=0cm 	% Permet d'enlever le trait vertical entre chaque colonne
%\renewcommand\tabenumsep{0.5cm\hskiplem}	% Ecarte de 0.5cm les items de tabenum, sinon ils sont trop collés...


\newlength{\tagmarginsep} 		% On crée une longueur
\setlength{\tagmarginsep}{0cm}		% Permet de régler la position d'un tag pour une équation mettre avant l'équation : \setlength{\tagmarginsep}{x cm} pour que le tag se rapproche de l'équation
\everydisplay{\displayindent=\tagmarginsep \displaywidth=\dimexpr\linewidth-2\tagmarginsep}	% Va avec ce qu'il y a juste avant


%-------------------- Définition des couleurs-----------------------------------------------------------------------------------------

\definecolor{violetfonce}{RGB}{126,62,142}
\definecolor{magentafonce}{RGB}{128, 0, 128}
\newrgbcolor{magentafonce}{.5 0 .5}			% Une autre façon de définir une couleur, permet surtout d'utiliser \magenta
\definecolor{violet1}{cmyk}{0.5,1,0,0}
\definecolor{violet2}{cmyk}{0.24998, 1, 0, 0}
\definecolor{violet3}{cmyk}{0, 1, 0.25, 0}



\newcommand\tred{\textcolor{red}}
\newcommand\tredf{\textcolor{rougefonce}}
\newcommand\tdred{\textcolor{deepred}}
\newcommand\tblue{\textcolor{blue}}
\newcommand\tbluef{\textcolor{bleufonce}}
\newcommand\tgreen{\textcolor{vertfonce}}
\newcommand\torange{\textcolor{orange}}
\newcommand\tmagenta{\textcolor{magentafonce}}


%------------------ Changement paramètres des parties, sous parties---------------------------------------------------------


\sectionfont{\color{violet1}}			% Utilisation du package sectsty
\subsectionfont{\color{violet2}}
\subsubsectionfont{\color{violet3}}


%--------------------------------- Raccourcis probabilités--------------------------------------------------------------------------------
\newcommand{\Prob}{\mathbb{P}}
\newcommand{\E}{\mathbb{E}}
\newcommand{\normale}[2]{\mathcal{N}\left(#1;#2 \right)}	% Pour la loi Normale : 	$\normale{\mu}{\sigma^2}$
\newcommand{\Unif}[1]{\mathcal{U}\left(#1\right)}			% Pour la loi uniforme :	$\Unif{[a;b]}$
\newcommand{\Ber}[1]{\mathcal{B}\left(#1\right)}			% Pour la loi de Bernoulli : $\Ber{p}$
\newcommand{\Bino}[2]{\mathcal{B}\left(#1;#2\right)}		% Pour la loi Binomiale :	$\Bino{n}{p}$
\newcommand*{\binomial}[2]{% 
  \ensuremath{ 
    \begin{pmatrix} 
      #2\\ 
      #1
    \end{pmatrix}}}				% Pour le coefficient binomial : $\binom{k}{n}$
\renewcommand{\binom}{\binomial} 	% La commande \binom existait déjà mais n'était pas pratique



%----------------------------------- Opérateurs---------------------------------------------------------------------------------------------------

\newcommand{\barre}[1]{\overline{#1}}  									% Barre au dessus du texte, pratique pour le conjugué des complexes
\newcommandx{\Sum}[3][3=\white]{{\displaystyle\sum\limits_{\substack{#1 #3}}^{#2}}}		% $\Sum{k=0}{n}[<option>]$ où [<option>] est placée à la suite de k=0 si l'on tape de façon classique, et en-dessous si l'on met [\\ k\neq k_0]
															% On peut juste taper : $\Sum{k=0}{n}$ et cela fonctionne
															% \displaystyle impose de gérer les tailles comme en mode mathématique : ici on aura donc toujours de grande somme que ce soit en math comme dans le texte
\newcommandx{\Prod}[3][3=\white]{{\displaystyle\prod\limits_{\substack{#1 #3}}^{#2}}}		% $\Prod{k=0}{n}[<option>]$	idem pour l'option
\newcommand{\Int}[4]{{\displaystyle\int_{#1}^{#2} #3  \d #4}}   					% $\Int{0}{\pi}{\cos(t)}{t}$
\newcommandx{\Lim}[3][3=\white]{\lim\limits_{\substack{#1 \to #2 #3}}}		% $\Lim{x}{0}[<option>]$	idem pour l'option


%------------------------------------ Délimiteurs--------------------------------------------------------------------------------------------------
%\newcommand{\abs}[1]{\left\lvert #1 \right\rvert}			% Valeur absolue | |
%\newcommand{\norm}[1]{\left\lVert #1 \right\rVert}			% Doubles barres verticales : norme d'un vecteur || ||

%--------------------------------------- Vecteurs----------------------------------------------------------------------------------------------------
 
\newcommand{\Oij}{\left(O;\vec{\imath},\vec{\jmath}\right)}
\newcommand{\Oijk}{\left(O;\vec{\imath},\vec{\jmath},\vec{k}\right)}
\newcommand{\Ouv}{\left(O;\vec{u},\vec{v}\right)}
\newcommand{\Ouvw}{\left(O;\vec{u},\vec{v},\vec{w}\right)}

%------------------------------------ Systèmes d'équations-------------------------------------------------------------------------------------------
% Remarque \syst et \syste ne font pas d'alignement mais ils sont bien, si l'on veut de l'alignement, on utilisera \systeme du package "systeme"
\newcommand*{\syst}[2]{%			% Systèmes 2 éq. à ? inconnues : 	\syst{ax+by+c=0}{a'x+b'y+c'=0}
  \left\{\begin{array}{ll} 
        #1\\ 
        #2 
      \end{array}  
    \right. }
 
\newcommand*{\syste}[3]{% 			% Systèmes 3 éq. à ? inconnues : 	\syste{ax+by+c=0}{a'x+b'y+c'=0}{a''x+b''y+c''=0}
  \left\{\begin{array}{ll} 
        #1\\ 
        #2\\
        #3 
      \end{array}  
    \right.}



%------------------------------------------ Autre----------------------------------------------------------------------------------------------------------
\newcommand{\boite}[2]{\noindent\parbox[]{#1 cm}{#2} }  									% Boite autre façon de faire minipage en gros


%---------------------------- En-têtes et bas de pages----------------------------------------------------------------------------
% A chaque fois : *head* pour le haut et *foot* pour le bas
% L = left, C = center, R = right 
% L = toutes les pages, LE = toutes les pages paires, LO = toutes les pages impaires, idem CE, CO, ...


\newcommand{\fcours}[1]{								% f pour fancy				   
	\pagestyle{fancy}									% Efface le style qui a pu être mis avant
	\renewcommand{\headrulewidth}{0pt} 
	\renewcommand{\footrulewidth}{0.5pt} 
	%\fancyhead[L]{}									
	%\fancyhead[C]{}
	%\fancyhead[R]{}									
	\fancyfoot[L]{\vspace{-1mm}\footnotesize \textit{#1}} 			% Classe : <Seconde> 
	\fancyfoot[C]{\vspace{-1mm}\footnotesize \thepage/\pageref{LastPage}}
	\fancyfoot[R]{\vspace{-1.8mm}\textit{\footnotesize C. Marguerite}}	% N'apparaîtra pas si on met les marges avec moins de bottom=1.3cm dans le package geometry
 }												% Fin de la commande
					%% Mettre au début du document : \fcours{Seconde} par exemple
					%% \usepackage[a4paper,top=0.4cm,bottom=1cm,left=0.4cm,right=0.4cm,marginparwidth=0.5cm,footskip=0.6cm]{geometry} 




%---------------------------- Réglages chapitres + table des matières------------------------------------------------------
\titleformat{\part}[display]
{\centering\normalfont\Huge\bfseries}
{\partname}
{20pt}
{\rule{6cm}{1pt}\raisebox{-2.1pt}[10pt][10pt]{\quad\decofourleft\decoone\decofourright\quad}\rule{6cm}{1pt} ~\\\vspace{20pt}}

\makeatletter
\renewcommand\chapter{\global\@topnum\z@
                    \@afterindentfalse
                    \secdef\@chapter\@schapter}
\def\@topnewpage[#1]{#1}
\makeatother										% Pour avoir le chapitre sur une demie page

\titleformat{\chapter}{\centering\normalfont\Large\bfseries\scshape\color{magentafonce}}{}{1em}{}
\titlespacing*{\chapter}{0pt}{-18pt plus 0pt minus 2pt}{0pt plus 2pt minus 2pt}		% Réglages espacements (à gauche, au-dessus, en-dessous) de la commande \chapter

\makeatletter
\newcommand{\twocolumntoc}{		
  \section*{\huge\centering\contentsname
    \@mkboth{%
      \MakeUppercase\contentsname}{\MakeUppercase\contentsname}}%
  \@starttoc{toc}%
}
\makeatother										% Permet de mettre la table des matières en twocolumn 
												% Mettre après \begin{doc..} \twocolumntoc \clearpage 		% Remplace \tableofcontents par \twocolumntoc en twocolumn

\makeatletter
\renewcommand{\thesection}{\@arabic\c@section}
\makeatother										% Enlève la numérotation des chapitres devant les sections

\setcounter{tocdepth}{3} 								% Permet d'afficher dans la table des matières suivant une certaine profondeur : -1=part, 0=chapter, 1=section, ...
\setcounter{secnumdepth}{3}							% Pour avoir la numérotation des subsubsections avec les classes book et report

\newcommand{\nocontentsline}[3]{}
\let\oldaddcontentsline\addcontentsline
\newcommand{\tocless}[2]{\let\addcontentsline\nocontentsline #1{#2}\let\addcontentsline\oldaddcontentsline}	% Pour retirer une section de la table des matières : \tocless\section{...}





\setstretch{1}
\begin{document}
\fcours{Concours général}


\chapter{Concours général des lycées - Session 2020}

\section{Problème 1 : Nombres pointus}

\n Soit $n$ un entier naturel non nul. On dit que $n$ est pointu si $n$ admet au plus un facteur premier ou bien si, en notant $p$ et $q$ les deux plus grands facteurs premiers de $n$, avec $p>q$, l'inégalité $p \geq 2 q$ est vérifiée.\\
Par exemple, 1 est pointu, car il n'a aucun facteur premier. De même, 25 est pointu, car il n'a qu'un seul facteur premier, et 147 est pointu, car $147=3 \times 7^{2}$ et $7 \geq 2 \times 3$. Au contraire, 105 n'est pas pointu, puisque $105=3 \times 5 \times 7$ et $7<2 \times 5$.\\
Dans ce problème, on cherche à démontrer qu'il existe des suites arbitrairement longues d'entiers consécutifs pointus. Plus précisément, on souhaite démontrer la propriété $\mathscr{P}$ suivante :

\begin{center}
Pour tout entier $m \geq 1$, il existe un entier $n \geq 0$ tel que les nombres $n+1, n+2, \ldots, n+m$ soient tous pointus.
\end{center}


\subsection{Quelques exemples}

\begin{enumerate}
	\item Le nombre 2020 est-il pointu ?
	\item Quel est le plus petit entier naturel non nul qui ne soit pas pointu?
	\item Quel est le plus petit nombre pointu possédant au moins quatre facteurs premiers distincts?
	\item Démontrer qu'il existe une infinité de nombres pointus.
	\item Démontrer qu'il existe une infinité d'entiers naturels non nuls qui ne sont pas pointus.
	\item Établir la liste des nombres pointus entre 1 et 20 inclus. Quelle est la longueur maximale d'une suite de nombres pointus consécutifs entre 1 et 20 ?
\end{enumerate}

\subsection{Peu de grands nombres premiers}


\n On pose $0 !=1$, et $\ell !=1 \times 2 \times \cdots \times \ell=\ell(\ell-1) !$ pour tout entier $\ell \geq 1$. Soient alors $k$ et $n$ deux entiers naturels tels que $k \leq n$. On s'intéresse à la fraction
$$\frac{n !}{k !(n-k) !}$$
que l'on note $F_{n, k}$.

\begin{enumerate}\setcounter{enumi}{6}
	\item \begin{enumerate}
		\item Calculer les valeurs des nombres $F_{3,1}$ et $F_{9,4}$.
		\item Démontrer que, si $k=0$ ou $k=n$, alors $F_{n, k}=1$.
		\item Démontrer que, si $1 \leq k \leq n-1$, alors $F_{n, k}=F_{n-1, k}+F_{n-1, k-1}$
		\item En déduire que, pour tout entier naturel $n$ et pour tout entier naturel $k \leq n, F_{n, k}$ est un entier naturel non nul inférieur ou égal à $2^{n}$.
	\end{enumerate} 
\end{enumerate}

\n Dans cette question et dans les parties qui suivent, pour tout entier naturel $n$, on note $\mathbb{P}_{n}$ l'ensemble des nombres premiers $p$ tels que $n+1 \leq p \leq 2 n$, et on note $\pi_{n}$ le nombre d'éléments de $\mathbb{P}_{n}$.

\begin{enumerate}\setcounter{enumi}{7}
	\item \begin{enumerate}
		\item Démontrer que, pour tout nombre premier $p$ appartenant à l'ensemble $\mathbb{P}_{n}$, l'entier $F_{2 n, n}$ est divisible par $p$.
		\item Démontrer que, si $a$, $b$ et $c$ sont des entiers naturels non nuls tels que $b$ et $c$ sont premiers entre eux et divisent $a$, alors l'entier $b c$ divise $a$ lui aussi.
		\item Soit $d$ le produit de tous les éléments de $\mathbb{P}_{n}$. Démontrer que l'entier $F_{2 n, n}$ est divisible par $d$.
		\item En déduire que $n^{\pi_{n}} \leq 2^{2 n}$.
	\end{enumerate}
\end{enumerate}


\subsection{Des ensembles à peu d'éléments}

\n Soient $b$ un entier naturel non nul, et $f_{b}:] 0 ;+\infty[\rightarrow \mathbb{R}$ la fonction définie par
$$f_{b}(x)=\frac{(\ln (x))^{b}}{x} \q \text { pour tout } x>0$$

\begin{enumerate}\setcounter{enumi}{8}
	\item \begin{enumerate}
		\item Soient $x$ un nombre réel tel que $x>1$, et soit $y=\ln (x) / b$. Démontrer que : 
		$$f_{b}(x)=\frac{b^{b}}{(\exp (y) / y)^{b}}$$
		\item En déduire que $\Lim{x}{+\infty} f_{b}(x)=0$.
	\end{enumerate}	 
\end{enumerate}

\n Dans la suite de cette partie, on considérera deux entiers naturels non nuls $\ell$ et $n$ fixés, tels que $n \geq 2$.

\begin{enumerate}\setcounter{enumi}{9}
	\item On note $A_{\ell}$ l'ensemble des entiers naturels non nuls dont tous les facteurs premiers sont inférieurs ou égaux à $2^{\ell}$.
	\begin{enumerate}
		\item Parmi les entiers compris entre 1 et 30 inclus, lesquels appartiennent à l'ensemble $A_{2}$ ?
		\item On note $p_{1}, p_{2}, \ldots, p_{k}$ les nombres premiers inférieurs ou égaux à $2^{\ell}$, numérotés dans l'ordre croissant. Soit alors $a$ un élément de $A_{\ell}$ qui soit inférieur ou égal à $n$.\\
		Soient $\alpha_{1}, \alpha_{2}, \ldots, \alpha_{k}$ les entiers naturels tels que :
		$$a=p_{1}^{\alpha_{1}} \times p_{2}^{\alpha_{2}} \times \cdots \times p_{k}^{\alpha_{k}}$$
		Démontrer que les entiers naturels $\alpha_{1}, \alpha_{2}, \ldots, \alpha_{k}$ sont tous inférieurs ou égaux à $\ln (n) / \ln (2)$.
		\item En déduire que $A_{\ell}$ compte au plus
		$$\left(\frac{\ln (n)}{\ln (2)}+1\right)^{2^{\ell}}$$
		éléments inférieurs ou égaux à $n$.
		\item On pose $b=2^{\ell}$. Démontrer que $A_{\ell}$ compte au plus
		$$\frac{2}{(\ln (2))^{b}} \times f_{b}(2 n) \times n$$
		éléments inférieurs ou égaux à $n$.
	\end{enumerate}
	\item Soit $k$ un entier naturel non nul. On rappelle que l'ensemble $\mathbb{P}_{2^k}$ désigne l'ensemble des nombres premiers $p$ tels que $2^{k}+1 \leq p \leq 2 \times 2^{k}$, et que $\pi_{2^{k}}$ est le nombre d'éléments de $\mathbb{P}_{2^{k}}$.
	\begin{enumerate}
		\item Expliciter les ensembles $\mathbb{P}_{2^{1}}$ et $\mathbb{P}_{2^{2}}$.
	\end{enumerate}
	On note $B_{k, 1}$ l'ensemble des entiers naturels non nuls ayant au moins un facteur premier dans l'ensemble $\mathbb{P}_{2^{k}}$, et on note $B_{k, 2}$ l'ensemble des entiers naturels non nuls ayant au moins deux facteurs premiers distincts dans l'ensemble $\mathbb{P}_{2^{k}}$.
	\begin{enumerate}\setcounter{enumii}{1}
		\item Quel est le plus petit élément de l'ensemble $B_{2,2}$ ?
		\item Soit $p$ un nombre premier. Démontrer que le nombre d'entiers naturels non nuls inférieurs ou égaux à $n$ et multiples de $p$ est au plus de $n / p$.
		\item Démontrer que l'ensemble $B_{k, 1}$ compte au plus $n \times \pi_{2^{k}} / 2^{k}$ éléments inférieurs ou égaux à $n$.
		\item En déduire que l'ensemble $B_{k, 1}$ compte au plus $2 n / k$ éléments inférieurs ou égaux à $n$.
		\item Soit $p_{1}$ et $p_{2}$ deux nombres premiers distincts. Démontrer que le nombre d'entiers naturels non nuls inférieurs ou égaux à $n$ et multiples de $p_{1}$ et $p_{2}$ est au plus de $n/\left(p_{1} \times p_{2}\right)$.
		\item En déduire que l'ensemble $B_{k, 2}$ compte au plus $4 n / k^{2}$ éléments inférieurs ou égaux à $n$.
	\end{enumerate}
	\item Pour tout entier $k \geq 2$, on note également $C_{k}$ l'ensemble $\left(B_{k-1,1} \cap B_{k, 1}\right) \cup B_{k, 2}$.
	\begin{enumerate}
		\item Quel est le plus petit élément de l'ensemble $C_{2}$ ?
		\item Démontrer que l'ensemble $C_{k}$ compte au plus :
		$$\frac{4 n}{k^{2}}+\frac{4 n}{k(k-1)}$$
		éléments inférieurs ou égaux à $n$.
	\end{enumerate}
	\item On fixe maintenant un entier $\ell \geq 2$, et on note $D_{\ell}$ l'ensemble des entiers qui appartiennent à au moins l'un des ensembles $C_{k}$ tels que $k \geq \ell$.
	\begin{enumerate}
		\item Démontrer que tout élément de $\{1, \ldots, n\}$ appartenant à l'ensemble $D_{\ell}$ appartient à l'un des ensembles $C_{\ell}, C_{\ell+1}, \ldots, C_{n}$
		\item Démontrer, pour tout entier $k \geq 2$, que :
		$$\frac{1}{k^{2}}+\frac{1}{k(k-1)} \leq \frac{2}{k-1}-\frac{2}{k}$$
		\item En déduire que $D_{\ell}$ compte au plus $8 n /(\ell-1)$ éléments inférieurs ou égaux à $n$.
	\end{enumerate}
\end{enumerate}


\subsection{Beaucoup de nombres pointus}

\n Dans cette partie, on considère un entier naturel non nul $m$ fixé.

\begin{enumerate}\setcounter{enumi}{13}
	\item Démontrer qu'il existe un entier $\ell \geq 2$ tel que, pour tout entier $n \geq 1$, l'ensemble $D_{\ell}$ compte au plus $n$ éléments inférieurs ou égaux à $3 mn$.
\end{enumerate}

\n Dans la suite, on considère également un tel entier $\ell$.

\begin{enumerate}\setcounter{enumi}{14}
	\item Soit $a$ un entier naturel non nul. Démontrer que, si $a$ n'est pas pointu, alors $a$ appartient à $A_{\ell}$ ou à $D_{\ell}$.
	\item Démontrer qu'il existe un entier $n \geq 1$ tel que l'ensemble $A_{\ell}$ compte au plus $n$ éléments inférieurs ou égaux à $3mn$.
	\item En déduire qu'il existe un entier naturel $k$ inférieur ou égal à $2 n$ et tel que chacun des entiers $m k+1$, $m k+2, \ldots, m(k+1)$ soit pointu.
	\item Démontrer la propriété $\mathscr{P}$ introduite en préambule de l'énoncé.
\end{enumerate}


\section{Problème 2 : Un nombre explosif}


\n Si $\alpha$ est un réel strictement positif, on définit la suite $\left(x_{n}\right)$ par $x_{1}=\alpha$ et 
$$x_{n+1}=\left(1+\frac{1}{x_{n}}\right)^{n} \q \text{ pour tout entier } n \geq 1,$$
que l'on appelle \textit{suite associée à $\alpha$.}\\
On dit que le nombre $\alpha$ est \textit{explosif} si la suite $\left(x_{n}\right)$ associée à $\alpha$ diverge vers $+\infty$.\\
Le but du problème est de prouver qu'il existe un unique nombre explosif.\\

\n Pour tout entier $n \geq 1$, on désigne par $f_{n}$ la fonction définie sur $\mathbb{R}_{+}^{*}$ par :
$$f_{n}(x)=\left(1+\frac{1}{x}\right)^{n} \q\text{ pour tout réel } x>0$$
Ainsi, pour tout entier $n \geq 1$, on a :
$$x_{n+1}=f_{n}\left(x_{n}\right)$$
Dans tout le problème, pour tout réel $t$, l'exponentielle de $t$ sera notée $\exp(t)$.

\subsection{Un encadrement de $f_{n}(x)$}

\n Soit $x$ un réel strictement positif.

\begin{enumerate}
	\item Démontrer que $\frac{1}{x+1} \leq \Int{x}{x+1}{\frac{1}{t}}{t}\leq \frac{1}{x}$.
	\item En déduire que $\frac{1}{x+1} \leq \ln \left(1+\frac{1}{x}\right) \leq \frac{1}{x}$.
	\item En déduire que, pour tout entier $n \geq 1$, on a :
	$$\exp (n /(x+1)) \leq f_{n}(x) \leq \exp (n / x)$$
\end{enumerate}


\subsection{Un critère d'explosivité}

\begin{enumerate}\setcounter{enumi}{3}
	\item Soit $\alpha$ un réel strictement positif et $\left(x_{n}\right)$ la suite associée à $\alpha$. La suite $\left(x_{n}\right)$ peut-elle être majorée ? Peut-elle être convergente ?
\end{enumerate}

\n Pour tout entier $n \geq 1$, on pose $y_{n}=\frac{n+1}{\ln (n+2)}-1$.

\begin{enumerate}\setcounter{enumi}{4}
	\item \begin{enumerate}
		\item Étudier les variations de la fonction $g$ définie sur $[2 ;+\infty[$ par
		$$g(x)=\frac{x}{\ln (x+1)} \q\text { pour tout réel } x \geq 2$$
		\item En déduire que, pour tout entier $n \geq 8$, on a $y_{n}>\e$.
		\item Démontrer que, pour tout entier $n \geq 2$, on a $y_{n}>1$ et $y_{n-1} \leq \frac{n}{\ln \left(y_{n}\right)}$.
	\end{enumerate}	 
	\item Soit $\alpha$ un réel strictement positif et $\left(x_{n}\right)$ la suite qui lui est associée.
	\begin{enumerate}
		\item Démontrer que, pour tout entier $n \geq 8$, si $x_{n}<y_{n-1}$ alors $x_{n+1}>n+1$ et $x_{n+2}<\e$.
		\item En déduire que, si $\alpha$ est explosif, alors $x_{n} \geq y_{n-1}$ pour tout entier $n \geq 8$. 
		\item Démontrer que, si $\alpha$ est explosif, alors $x_{n} \leq \frac{n}{\ln \left(y_{n}\right)}$ pour tout entier $n \geq 8$.
		\item Démontrer que $\alpha$ est explosif si et seulement si $y_{n-1} \leq x_{n} \leq \frac{n}{\ln \left(y_{n}\right)}$ pour tout entier $n \geq 8$.
	\end{enumerate}
	\item Soient $\alpha, \beta$ et $\gamma$ des réels strictement positifs tels que $\alpha \leq \beta \leq \gamma$. Démontrer que, si $\alpha$ et $\gamma$ sont explosifs, alors $\beta$ est explosif.
\end{enumerate}


\subsection{Un unique nombre explosif}

\n Soient $r$ et $s$ deux fonctions définies sur $\mathbb{R}_{+}^{*}$ et telles que $s(x)>0$ pour tout réel $x>0$. On désigne par $r \circ s$ la fonction définie sur $\mathbb{R}_{+}^{*}$ par :
$$r \circ s(x)=r(s(x)) \q\text { pour tout réel } x>0$$
On pose alors $h_{1}=f_{1}$ et 
$$h_{n}=f_{n} \circ h_{n-1} \q\text { pour tout entier } n \geq 2$$
De plus, on note $h_{n}\left(\mathbb{R}_{+}^{*}\right)$ l'ensemble des réels $h_{n}(x)$ lorsque $x$ décrit $\mathbb{R}_{+}^{*}$.\\

\n On admettra enfin que, pour tout entier $n \geq 1$, la fonction $h_{n}$ est continue sur $\mathbb{R}_{+}^{*}$.

\begin{enumerate}\setcounter{enumi}{7}
	\item \begin{enumerate}
		\item Expliciter l'expression de $h_{2}(x)$ pour tout réel $x>0$.
		\item Soit $\alpha$ un réel strictement positif et $\left(x_{n}\right)$ sa suite associée.\\
		Pour tout entier $n \geq 1$, exprimer $x_{n+1}$ en fonction de $h_{n}$ et $\alpha$.
	\end{enumerate}	 
	\item \begin{enumerate}
		\item Démontrer que, pour tout entier $n \geq 1$, la fonction $h_{n}$ est strictement monotone sur $\mathbb{R}_{+}^{*}$.
		\item À l'aide de la calculatrice et sans plus de justification, déterminer des valeurs approchées à $10^{-2}$ près des réels $u$ et $v$ tels que $]u ; v[=h_{8}\left(\mathbb{R}_{+}^{*}\right)$.\\
		Vérifier que l'on a $u<\e$ et $9<v$.
		\item Démontrer que, pour tout entier $n \geq 9$, on a $[\e ; n] \subset h_{n-1}\left(\mathbb{R}_{+}^{*}\right)$.
		\item En déduire que, pour tout entier $n \geq 9$, on a $\left[y_{n-1} ; n / \ln \left(y_{n}\right)\right] \subset h_{n-1}\left(\mathbb{R}_{+}^{*}\right)$.
	\end{enumerate}
\end{enumerate}

\n Pour tout entier $n \geq 9$, on pose $I_{n}=\left[y_{n-1} ; n / \ln \left(y_{n}\right)\right]$ et on note $J_{n}$ l'ensemble des réels strictement positifs $x$ tels que $h_{n-1}(x) \in I_{n}$.

\begin{enumerate}\setcounter{enumi}{9}
	\item \begin{enumerate}
		\item Démontrer que, pour tout entier $n \geq 9$, il existe des réels $a_{n}$ et $b_{n}$ tels que $J_{n}=\left[a_{n} ; b_{n}\right]$.
		\item Démontrer que, pour tout entier $n \geq 9$, on a $J_{n+1} \subset J_{n}$.
		\item En déduire que les suites $\left(a_{n}\right)$ et $\left(b_{n}\right)$ convergent.
	\end{enumerate}
\end{enumerate} 

\n Dans toute la suite, on note $\alpha$ et $\beta$ les limites respectives des suites $\left(a_{n}\right)$ ) et $\left(b_{n}\right)$.\\
De plus, on désigne par $\left(\alpha_{n}\right)$ et $\left(\beta_{n}\right)$ les suites respectivement associées à $\alpha$ et $\beta$. On rappelle que cela signifie que $\alpha_{1}=\alpha$ et $\beta_{1}=\beta$ et que, pour tout entier $n \geq 1$ :
$$\alpha_{n+1}=f_{n}\left(\alpha_{n}\right) \q\text{ et }\q \beta_{n+1}=f_{n}\left(\beta_{n}\right)$$

\begin{enumerate}\setcounter{enumi}{10}
	\item \begin{enumerate}
		\item Démontrer que $\alpha$ et $\beta$ sont explosifs.
		\item Démontrer qu'un réel $x>0$ est explosif si et seulement si $\alpha \leq x \leq \beta$.
	\end{enumerate}
	\item \begin{enumerate}
		\item Démontrer que, pour tout entier $n \geq 9$ et tout réel $x \in I_{n}$, on a :
		$$\left|f_{n}'(x)\right| \geq\left|f_{n}'\left(\frac{n}{\ln \left(y_{n}\right)}\right)\right|$$
		\item En déduire que, pour tout entier $n \geq 9$ et tout réel $x \in I_{n}$, on a :
		$$\left|f_{n}'(x)\right| \geq \frac{y_{n}\left(\ln \left(y_{n}\right)\right)^{2}}{n} \exp \left(-\frac{1+\ln \left(y_{n}\right)}{n+\ln \left(y_{n}\right)} \ln \left(y_{n}\right)\right)$$
		\item Démontrer que $\Lim{n}{+\infty}\frac{y_{n} \ln (n)}{n}=1$ et que $\Lim{n}{+\infty}\frac{\ln \left(y_{n}\right)}{\ln (n)}=1$.
		\item En déduire qu'il existe un entier $N \geq 9$ tel que, pour tout entier $n \geq N$ et tout réel $x \in I_{n}$, on a :
		$$\left|f_{n}'(x)\right| \geq 2$$
	\end{enumerate}
	\item Dans cette question, on désigne donc par $N$ un tel entier.
	\begin{enumerate}
		\item Déduire du \textbf{12.d.} que, pour tout entier $n \geq N$ et pour tous réels $x$ et $y$ dans $I_{n}$, on a :
		$$\left|f_{n}(x)-f_{n}(y)\right| \geq 2|x-y|$$
		\item Démontrer que, si $\alpha<\beta$, alors il existe un réel strictement positif $C$ tel que :
		$$\left|\alpha_{n}-\beta_{n}\right| \geq C \times 2^{n} \q\text { pour tout entier } n \geq N$$
	\end{enumerate}
	\item \begin{enumerate}
		\item Démontrer qu'il n'existe qu'un seul nombre explosif.
		\item À l'aide de la calculatrice et en indiquant l'algorithme utilisé, donner une valeur approchée de $\alpha$ à $10^{-2}$ près.
	\end{enumerate}
\end{enumerate}


\end{document}